% Part: computability
% Chapter: computability-theory
% Section: normal-form

\documentclass[../../../include/open-logic-section]{subfiles}

\begin{document}

\olfileid{cmp}{thy}{nfm}
\olsection{నియత రూప సిద్ధాంతం}

గణనీయ ప్రమేయాల నిర్వచనాలను మనం వర్ణించగలమనీ, ఒక ప్రమేయాన్ని ఏదో
ఇన్‌పుట్‌పై గణించిన పూర్తి లేఖనానికి ప్రతిపాదించిన వివరణ సరైనదో కాదో
పరీక్షించగలమనీ అనుకుందాం. అప్పుడు గణనా నమూనాతో సంబంధం లేకుండా, ఏ
గణనీయ ప్రమేయం~$f$ యొక్క ఏ ఇన్‌పుట్~$x$ వద్ద విలువనైనా కింది విధంగా
నిర్ణయించవచ్చని సహజంగా అనిపిస్తుంది:
\begin{enumerate}
  \item గణనా లేఖనాల సాధ్యమైన వివరణలన్నింటిలో వెతకాలి.
  \item ఇచ్చిన లేఖనం $f(x)$ గణన యొక్క లేఖనమేనా అని పరీక్షించాలి.
  \item సరైన లేఖనం దొరికితే దాని నుంచి $f(x)$ విలువను వెలికితీయాలి.
\end{enumerate}
ఇది నిజమేనని చెప్పేదే క్లీనీ నియత రూప సిద్ధాంతం.

\begin{thm}[క్లీనీ నియత రూప సిద్ధాంతం]
\ollabel{thm:normal-form}
కింది ధర్మం గల ఆదిమ పునరావృత్త సంబంధం~$T(e,x,s)$, ఆదిమ పునరావృత్త
ప్రమేయం~$U(s)$ ఉన్నాయి: $f$ ఏదైనా పాక్షిక గణనీయ ప్రమేయమైతే, ఏదో
$e$కు ప్రతి~$x$కూ
\[
f(x) \simeq U(\umin{s}{T(e, x, s)})
\]
అవుతుంది.
\end{thm}

\begin{proof}[నిరూపణ రూపురేఖ]
ఏ గణనా నమూనాకైనా, గణనీయ ప్రమేయం~$f$ యొక్క వివరణను కచ్చితంగా
నిర్వచించి, ఆ వివరణను సహజ సంఖ్య~$e$తో సంకేతీకరించవచ్చు. అలాగే
ఇన్‌పుట్~$x$పై సూచిక~$e$ గల ప్రమేయాన్ని గణించే ప్రక్రియను లిఖించే
``గణనా అనుక్రమం'' అనే భావనను కచ్చితంగా నిర్వచించవచ్చు. అటువంటి గణనా
అనుక్రమాన్ని కూడా సంఖ్య~$s$తో సంకేతీకరించవచ్చు. ఇది కింది రెండూ
గణనీయమయ్యేలా చేయవచ్చు:
\begin{enumerate}
  \item సూచిక~$e$ గల ప్రమేయాన్ని ఇన్‌పుట్~$x$పై గణించిన అనుక్రమాన్ని
  సంఖ్య~$s$ సంకేతీకరిస్తే, అప్పుడే మాత్రమే నెరవేరే సంబంధం $T(e,x,s)$;
  \item $s$తో సంకేతీకరించిన గణనా అనుక్రమాన్ని ఆ గణన తుది ఫలితానికి
  పంపే ప్రమేయం $U(s)$.
\end{enumerate}
నిజానికి సంబంధం~$T$, ప్రమేయం~$U$ రెండూ ఆదిమ పునరావృత్తాలు.
\end{proof}

\begin{explain}
నియత రూప సిద్ధాంతానికి కచ్చితమైన నిరూపణ ఇవ్వాలంటే, ఒక గణనా నమూనాను
ఎంచుకుని, గణనీయ ప్రమేయాల వివరణలకూ గణనా అనుక్రమాలకూ సంకేతీకరణను
వివరంగా నిర్వహించాలి; సంబంధం~$T$, ప్రమేయం~$U$ ఆదిమ పునరావృత్తాలని
నిర్ధారించాలి. చాలా అనువర్తనాలకు $T$, $U$ గణనీయాలని, $U$
సర్వనిర్వచితమని తెలిస్తే చాలు.

నియత రూప సిద్ధాంతపు నిరూపణను ఒక నినాదంగా గుర్తుంచుకోవడం బహుశా ఉత్తమం:
$\umin{s}{T(e,x,s)}$ అనేది ఇన్‌పుట్~$x$పై సూచిక~$e$ గల ప్రమేయపు
గణనా అనుక్రమం కోసం వెతుకుతుంది; అలాంటి అనుక్రమం దొరికితే $U$ దాని
అవుట్‌పుట్‌ను తిరిగి ఇస్తుంది.
\end{explain}

గణనా నమూనా పాక్షిక పునరావృత్త ప్రమేయాలైతే (క్లీనీ మొదట వాడింది ఇదే),
ఏ ప్రమేయపు నిర్వచనానికైనా అపరిమిత అన్వేషణను ఒక్కసారి మాత్రమే, అంటే
$\umin{y}{f(\vec x,y)}$ క్రియాకారిని ఒక్కసారి మాత్రమే వాడడం అవసరమని
ఇది చూపుతుంది. ఈ అర్థంలో, ప్రతి పాక్షిక పునరావృత్త ప్రమేయానికి ఒక
నియత రూపం ఉంటుందని ఇది చూపుతుంది.

$\cfind{0}$, $\cfind{1}$, $\cfind{2}$, \dots{} అనే లెక్కింపును
నిర్వచించడానికి $T$, $U$లను వాడవచ్చు. ఇకపై తగిన $T$, $U$లను
స్థిరపరచుకున్నామని ఊహించి,
\[
\cfind{e}(x) \simeq U(\umin{s}{T(e,x,s)})
\]
అనే సమీకరణాన్నే $\cfind{e}$ యొక్క \emph{నిర్వచనం}గా తీసుకుంటాం.

మరొక ఉపయోగకరమైన వాస్తవం ఇదిగో:

\begin{thm}
ప్రతి పాక్షిక గణనీయ ప్రమేయానికి అనంతంగా అనేక సూచికలు ఉంటాయి.
\end{thm}

ఇది కూడా అంతర్బోధాత్మకంగా స్పష్టమే. ఒక గణనీయ ప్రమేయపు ఏ వివరణనైనా
తీసుకుంటే, అదే ప్రమేయాన్ని (అదే ఇన్‌పుట్--అవుట్‌పుట్ జతలను) గణించే
మరొక వివరణను తయారు చేయవచ్చు. ఉదాహరణకు, ముందుగా గణనపై ప్రభావం లేని
పని ఏదైనా చేయించవచ్చు ($0=0$ అని పరీక్షించడం, $5$ వరకు లెక్కించడం
వంటివి). మార్చిన వివరణ సూచిక ఎల్లప్పుడూ అసలు సూచికకు భిన్నంగా
ఉంటుంది. కొద్దిగా భిన్నంగా గణించినప్పటికీ, రెండూ ఒకే ప్రమేయపు
సూచికలే.

\end{document}
